% fig_teaser.tex -- Figure 1 (in the title block, full width): the idea in three panels.
%   (a) the nearest-neighbour distance has no unit: two inputs at the same distance r from their nearest train image
%   (b) one transposition as the unit: the train image and its transposed copy; the distance tau between their descriptors; r measured in units of tau (same geometry as (a))
%   (c) real scores (figs/fig_hist_in64.tex, written by gen/mk.py): ImageNet-64 against iNaturalist, distance only and TRA
% Colours: navy = train images (kept at inference), white square = the transposed copy (used once), red = the input / OOD.
\begingroup
\tikzset{tz/.style={font=\figlab, line cap=round, line join=round}, lab/.style={font=\figlab, inner sep=1pt, align=center},
  tr/.style={circle, fill=navyc, inner sep=0, minimum size=4.2pt}, cp/.style={rectangle, draw=navyc, fill=white, line width=0.6pt, inner sep=0, minimum size=4.0pt},
  inp/.style={star, star points=5, star point ratio=2.2, fill=redc, inner sep=1.2pt}}
\newcommand{\paneltitle}[1]{{\fontsize{7.4}{9}\selectfont\bfseries #1}\\[0.6mm]}
\noindent
\begin{minipage}[t][47mm][t]{35mm}\centering
\paneltitle{(a) 距離には単位がない}
\begin{tikzpicture}[tz]
\useasboundingbox (0,-0.45) rectangle (3.5,3.55);
\coordinate (A) at (0.75,1.85); \coordinate (B) at (2.45,1.85);
\node[tr] at (A) {}; \node[tr] at (B) {};
\draw[black, line width=0.6pt] (A) -- node[pos=0.5, above left, inner sep=0.6pt, lab] {$r$} ($(A)+(40:0.55)$);
\draw[black, line width=0.6pt] (B) -- node[pos=0.5, above left, inner sep=0.6pt, lab] {$r$} ($(B)+(40:0.55)$);
\node[inp] at ($(A)+(40:0.55)$) {}; \node[inp] at ($(B)+(40:0.55)$) {};
\node[lab, text=navyc] at ($(A)+(0,-0.3)$) {$x_i$}; \node[lab, text=navyc] at ($(B)+(0,-0.3)$) {$x_j$};
\node[lab] at (1.75,0.75) {同じ距離 $r$.\\[1pt]近いのか, 遠いのか?};
\node[tr] at (0.3,0.05) {}; \node[lab, anchor=west] at (0.42,0.05) {train 画像};
\node[inp, inner sep=0.95pt] at (2.0,0.05) {}; \node[lab, anchor=west] at (2.12,0.05) {入力};
\end{tikzpicture}
\end{minipage}\hfill
\begin{minipage}[t][47mm][t]{66mm}\centering
\paneltitle{(b) 転置1つを長さの単位にする}
\begin{tikzpicture}[tz, >={Stealth[length=1.5mm, width=1.2mm]}]
\useasboundingbox (0,-0.45) rectangle (6.6,3.55);
% the train image and its transposed copy
\node[inner sep=0, anchor=south west, draw=navyc, line width=1pt] (im0) at (0.75,2.62) {\includegraphics[width=0.78cm]{figs/o0.png}};
\node[inner sep=0, anchor=south west, draw=black!55, line width=0.5pt] (im1) at (2.55,2.62) {\includegraphics[width=0.78cm]{figs/o5.png}};
\draw[black!45, densely dashed, line width=0.4pt] (im0.north west) -- (im0.south east);
\draw[->, black, line width=0.6pt] (im0.east) -- node[above, lab, inner sep=1.5pt] {転置} (im1.west);
\node[lab, anchor=east, text=navyc] at (0.68,3.01) {$x_i$};
\node[lab, anchor=west, align=left] at (3.5,3.01) {$\mathbf{T}x_i$：縦と横を入れ替える.\\記述子が $\tau$ だけ動く.};
% left: a large transposition distance, the input closer than tau
\coordinate (A) at (1.55,1.25);
\draw[navyc!55, densely dashed, line width=0.5pt] (A) circle (0.9);
\draw[navyc, line width=0.7pt] (A) -- node[pos=0.5, below, inner sep=1.8pt, lab, text=navyc] {$\tau$} ($(A)+(200:0.9)$);
\node[cp] at ($(A)+(200:0.9)$) {}; \node[tr] at (A) {};
\draw[black, line width=0.6pt] (A) -- node[pos=0.5, above left, inner sep=0.6pt, lab] {$r$} ($(A)+(40:0.55)$);
\node[inp] at ($(A)+(40:0.55)$) {};
\node[lab] at (1.55,-0.12) {$\tau/r>1$：分布内らしい};
% right: a small transposition distance, the input farther than tau
\coordinate (B) at (4.75,1.25);
\draw[navyc!55, densely dashed, line width=0.5pt] (B) circle (0.27);
\draw[navyc, line width=0.7pt] (B) -- ($(B)+(200:0.27)$);
\node[cp] at ($(B)+(200:0.27)$) {}; \node[tr] at (B) {};
\draw[black, line width=0.6pt] (B) -- ($(B)+(40:0.55)$);
\node[lab] at ($(B)+(64:0.66)$) {$r$};
\node[inp] at ($(B)+(40:0.55)$) {};
\node[lab, text=navyc] at ($(B)+(-0.42,-0.38)$) {$\tau$};
\node[lab] at (4.75,-0.12) {$\tau/r<1$：分布外らしい};
\node[cp] at (2.74,0.62) {}; \node[lab, anchor=west] at (2.86,0.62) {転置した train 画像};
\draw[navyc!55, densely dashed, line width=0.5pt] (2.62,0.3) -- (2.88,0.3); \node[lab, anchor=west] at (2.92,0.3) {半径 $\tau$};
\end{tikzpicture}
\end{minipage}\hfill
\begin{minipage}[t][47mm][t]{70mm}\centering
\paneltitle{(c) ImageNet-64 と iNaturalist のスコアの分布}
\vspace{0.5mm}
% written by gen/mk.py from data/fig_tra.json (stored test scores, noise seed 0)
\begin{tikzpicture}
\begin{groupplot}[group style={group size=2 by 1, horizontal sep=7mm}, width=38mm, height=33mm,
  tick label style={font=\figlab}, label style={font=\figlab}, title style={font=\figlab, yshift=-1.6mm, align=center}, xlabel={標準化したスコア},
  legend style={font=\figlab, draw=none, fill=none, legend columns=2, /tikz/every even column/.append style={column sep=3mm}}]
\nextgroupplot[title={距離のみ\\AUROC 79.6\%}, ylabel={密度}, legend to name=histlegA]
\addplot[const plot, navyc, line width=0.7pt] coordinates {(-7.50,0.0000) (-7.00,0.0040) (-6.50,0.0040) (-6.00,0.0000) (-5.50,0.0000) (-5.00,0.0000) (-4.50,0.0000) (-4.00,0.0040) (-3.50,0.0120) (-3.00,0.0320) (-2.50,0.0360) (-2.00,0.0440) (-1.50,0.1240) (-1.00,0.1880) (-0.50,0.3920) (0.00,0.5120) (0.50,0.4040) (1.00,0.2160) (1.50,0.0280) (2.00,0.0000) (2.50,0)};
\addplot[const plot, redc, densely dashed, line width=0.7pt] coordinates {(-7.50,0.0000) (-7.00,0.0000) (-6.50,0.0000) (-6.00,0.0200) (-5.50,0.0240) (-5.00,0.0200) (-4.50,0.0560) (-4.00,0.0800) (-3.50,0.0840) (-3.00,0.1520) (-2.50,0.1720) (-2.00,0.2600) (-1.50,0.1760) (-1.00,0.3480) (-0.50,0.2360) (0.00,0.2600) (0.50,0.0960) (1.00,0.0160) (1.50,0.0000) (2.00,0.0000) (2.50,0)};
\legend{分布内の画像, 分布外の画像}
\nextgroupplot[title={\method\\AUROC 92.8\%}]
\addplot[const plot, navyc, line width=0.7pt] coordinates {(-4.50,0.0000) (-4.00,0.0000) (-3.50,0.0080) (-3.00,0.0040) (-2.50,0.0520) (-2.00,0.0600) (-1.50,0.2120) (-1.00,0.3120) (-0.50,0.3360) (0.00,0.3680) (0.50,0.2640) (1.00,0.2840) (1.50,0.1000) (2.00,0.0000) (2.50,0)};
\addplot[const plot, redc, densely dashed, line width=0.7pt] coordinates {(-4.50,0.0000) (-4.00,0.0040) (-3.50,0.0520) (-3.00,0.2400) (-2.50,0.4680) (-2.00,0.5800) (-1.50,0.4480) (-1.00,0.1440) (-0.50,0.0520) (0.00,0.0080) (0.50,0.0040) (1.00,0.0000) (1.50,0.0000) (2.00,0.0000) (2.50,0)};
\end{groupplot}
\end{tikzpicture}\\[-0.5mm]
\pgfplotslegendfromname{histlegA}

\end{minipage}
\endgroup
